% ============================================================
%  REGISTRE DES RISQUES — Phase 2 : Planification (MAJ mensuelle)
%  Time'Eats — Cellule PMO
%  Projet : Refonte Web
% ============================================================
\documentclass[11pt,a4paper,oneside]{report}
\usepackage{../../style/consulting}
\usepackage{pdflscape}
\hypersetup{pdftitle={Registre des Risques — Refonte Web — Time'Eats PMO}}
\pagestyle{fancy}\fancyhf{}
\renewcommand{\headrulewidth}{0pt}
\fancyhead[L]{\begin{tikzpicture}[remember picture,overlay]
  \draw[cnNavy,line width=0.5pt](0,0)--(\textwidth,0);
\end{tikzpicture}\vspace{2pt}\timeeatslogo\hspace{4pt}\small\textcolor{cnSlate}{Time'Eats \textbf{·} Registre des Risques --- \textit{Refonte Web}}}
\fancyhead[R]{\small\textcolor{cnSlate}{\thepage}}
\fancyfoot[C]{\tiny\textcolor{cnSlate}{Document confidentiel --- Cellule PMO --- CESI MAALSI 2026}}
\begin{document}\small

\begin{center}
{\LARGE\bfseries\color{cnNavy} Registre des Risques}\\[4pt]
\textcolor{cnOrange}{\rule{10cm}{1pt}}\\[4pt]
{\large\color{cnSlate} Phase 2 --- Planification (document vivant, MAJ mensuelle)}
\end{center}

\vspace{4pt}
\renewcommand{\arraystretch}{1.4}
\begin{tabularx}{\textwidth}{L{3.5cm} Y L{3.5cm} Y}
\toprule
\rowcolor{cnNavy}\th{Projet} & \th{Refonte Web (REF-WEB-2026)} & \th{Dernière MAJ} & \th{19/05/2026} \\
\midrule
Chef de projet  & M. X & Version & 1.2 (post-COPIL inception) \\
\bottomrule
\end{tabularx}

\vspace{6pt}
\begin{finding}
\textbf{Grille de cotation :}
Probabilité 1=rare, 2=peu probable, 3=possible, 4=probable, 5=quasi-certain. \\
Impact 1=négligeable, 2=mineur, 3=modéré, 4=majeur, 5=critique. \\
Score = P $\times$ I. \textbf{Critique $\geq$ 15 / Élevé 8--14 / Modéré 4--7 / Faible $\leq$ 3.}
\end{finding}

\begin{landscape}
\renewcommand{\arraystretch}{1.4}
\begin{tabularx}{\linewidth}{C{0.5cm} L{3.5cm} C{0.7cm} C{0.7cm} C{0.8cm} L{2.8cm} L{4cm} C{2.2cm} C{1.6cm}}
\toprule
\rowcolor{cnNavy}
\th{\#} & \th{Description du risque} & \th{P} & \th{I} & \th{Score} & \th{Stratégie} & \th{Plan de mitigation} & \th{Propriétaire} & \th{Statut} \\
\midrule

R1 & Retard projet Office 365 / SSO non livré à temps pour l'espace client (J3)
   & 3 & 4 & \cellcolor{cnOrange!40}\textbf{12}
   & Réduire (repli technique)
   & Repli IdP local prévu (surcoût 8 K€). Synchronisation hebdo avec PMO Office 365.
   & Tech Lead (M. A) & \cellcolor{cnOrange!20}Ouvert \\

\rowcolor{cnIce}
R2 & Non-conformité RGAA 4.1 niveau AA à l'audit J4 (refonte UX insuffisante)
   & 3 & 4 & \cellcolor{cnOrange!40}\textbf{12}
   & Réduire
   & Pré-audit à mi-parcours (S30), revue accessibilité par UX experte, formation contributeurs CMS.
   & UX (Mme C) + RSSI & \cellcolor{cnOrange!20}Ouvert \\

R3 & Sous-effectif ESN Sopra Steria sur sprint critique (été, congés)
   & 3 & 3 & \cellcolor{cnOrange!40}\textbf{9}
   & Transférer (contractuel)
   & Clause contractuelle forfaitaire de continuité (BC-2026-0142). Backup nominatif identifié côté ESN.
   & CP (M. X) & \cellcolor{cnOrange!20}Ouvert \\

\rowcolor{cnIce}
R4 & Indisponibilité référent Marketing (Mme F) compromettant la validation des contenus
   & 2 & 4 & \cellcolor{cnOrange!40}\textbf{8}
   & Réduire
   & Désigner un référent suppléant Marketing (en cours). Créneaux récurrents mardis 14h--17h sanctuarisés.
   & DSI + Marketing & \cellcolor{cnOrange!20}Ouvert \\

R5 & Dépendance migration Azure West Europe non finalisée à J2 (socle technique)
   & 2 & 4 & \cellcolor{cnOrange!40}\textbf{8}
   & Réduire
   & Coordination tech weekly avec équipe Infra. Solution de repli : tenant Azure existant en isolation.
   & DevOps (M. D) & \cellcolor{cnOrange!20}Ouvert \\

R6 & Performance Web Vitals dégradées (LCP $>$ 2{,}5 s sur mobile) à la recette
   & 3 & 3 & \cellcolor{cnOrange!40}\textbf{9}
   & Réduire
   & Tests Lighthouse en continu (pipeline CI), budget perf $<$ 200 ko JS par page, CDN Azure Front Door.
   & Tech Lead & \cellcolor{cnOrange!20}Ouvert \\

\rowcolor{cnIce}
R7 & Vulnérabilité critique détectée au pentest OWASP ASVS 2 (J4)
   & 2 & 5 & \cellcolor{cnRed!40}\textbf{10}
   & Réduire
   & SAST en CI (CodeQL + Dependabot), threat modeling au démarrage, code review systématique par RSSI sur modules sensibles (auth, paiement).
   & RSSI (Mme E) & \cellcolor{cnOrange!20}Ouvert \\

\bottomrule
\end{tabularx}
\end{landscape}

\section*{1 --- Matrice de criticité (Heat Map)}

\begin{center}
\renewcommand{\arraystretch}{1.6}
\begin{tabular}{c|ccccc}
 & \multicolumn{5}{c}{\textbf{Impact $\rightarrow$}} \\
\textbf{Proba $\downarrow$} & 1 & 2 & 3 & 4 & 5 \\
\hline
5 & \cellcolor{cnOrange!40}5 & \cellcolor{cnOrange!40}10 & \cellcolor{cnRed!40}15 & \cellcolor{cnRed!40}20 & \cellcolor{cnRed!70}25 \\
4 & \cellcolor{cnGreen!30}4 & \cellcolor{cnOrange!40}8 & \cellcolor{cnOrange!40}12 & \cellcolor{cnRed!40}16 & \cellcolor{cnRed!40}20 \\
3 & \cellcolor{cnGreen!30}3 & \cellcolor{cnGreen!30}6 & \cellcolor{cnOrange!40}\textbf{R3,R6} & \cellcolor{cnOrange!40}\textbf{R1,R2} & \cellcolor{cnRed!40}15 \\
2 & \cellcolor{cnGreen!30}2 & \cellcolor{cnGreen!30}4 & \cellcolor{cnGreen!30}6 & \cellcolor{cnOrange!40}\textbf{R4,R5} & \cellcolor{cnOrange!40}\textbf{R7} \\
1 & \cellcolor{cnGreen!30}1 & \cellcolor{cnGreen!30}2 & \cellcolor{cnGreen!30}3 & \cellcolor{cnGreen!30}4 & \cellcolor{cnGreen!30}5 \\
\end{tabular}
\end{center}

\section*{2 --- Synthèse criticité}

\renewcommand{\arraystretch}{1.3}
\begin{tabularx}{\textwidth}{L{3cm} C{2.5cm} Y}
\toprule
\rowcolor{cnNavy}\th{Niveau} & \th{Nombre} & \th{Identifiants} \\
\midrule
\cellcolor{cnRed!20}\textbf{Critique} ($\geq$ 15)   & 0  & --- \\
\rowcolor{cnIce}
\cellcolor{cnOrange!20}\textbf{Élevé} (8--14)       & 7  & R1, R2, R3, R4, R5, R6, R7 \\
\cellcolor{cnGreen!20}\textbf{Modéré} (4--7)        & 0  & --- \\
\rowcolor{cnIce}
\cellcolor{cnGreen!20}\textbf{Faible} ($\leq$ 3)    & 0  & --- \\
\bottomrule
\end{tabularx}

\vspace{4pt}
\begin{alert}
\textbf{Profil de risque actuel :} 100\,\% des risques identifiés sont en zone \textbf{Élevée} --- aucun critique, mais une surveillance soutenue est requise. Revue mensuelle obligatoire en COPIL. Les risques R1, R2 et R7 font l'objet d'une fiche d'action dédiée.
\end{alert}

\begin{reco}
À chaque MAJ mensuelle, le CP :
\begin{enumerate}
  \item Réévalue P et I de chaque risque ouvert.
  \item Ajoute les nouveaux risques émergents (issus des CR de sprint, COPIL, audits).
  \item Clôture les risques dont la fenêtre temporelle est passée (ex.\ R3 après mi-août).
  \item Reporte les 3 risques les plus critiques dans le rapport d'avancement (P3-13).
\end{enumerate}
\end{reco}

\end{document}
